% === §3 Experimental Setup ===
\section{Experimental Setup}
\label{sec:experimental-setup}

\looseness=-1
We evaluate \sysname{} on four benchmark families with seven domain configurations, covering long-term conversational question answering, embodied household task completion, medical decision support, and professional legal and financial reasoning.
For each configuration, we run 20 evolution iterations, select the best memory program on validation, and report performance on a held-out test split.
We use GPT-5.4 Mini as the task agent and GPT-5.3-Codex as the coding agent in reflective code evolution~\citep{openai2026gpt54mini,openai2026gpt53codex}.
Additional dataset splits, hyperparameters, evaluation protocols, computational costs, and pseudocode are provided in Appendix~\ref{app:datasets}, Appendix~\ref{app:hyperparameters}, Appendix~\ref{app:eval-protocols}, Appendix~\ref{app:cost}, and Appendix~\ref{app:algorithm}.

\noindent\textbf{Benchmarks.}
We use \rev{LoCoMo~\citep{locomo}} to test long-term conversational memory, \rev{ALFWorld~\citep{alfworld}} to test memory use in embodied planning, \rev{HealthBench~\citep{healthbench}} to test rubric-graded health reasoning, and \rev{PRBench~\citep{prbench}} to test professional legal and financial reasoning.
Together, these benchmarks cover retrieval-heavy question answering, interactive decision making, and specialized reasoning under structured evaluation.
Dataset construction details are provided in Appendix~\ref{app:datasets}, and metric details are provided in Appendix~\ref{app:eval-protocols} for all seven benchmark configurations.

\noindent\textbf{Baselines.}
We compare \sysname{} with nine baselines, including a no-memory control and eight memory-based methods spanning three paradigms, as summarized in Table~\ref{tab:main-results}.
These baselines cover direct answering without external memory, retrieval-based memory over stored observations, self-evolving memories that distill reusable experience, and prompt-optimization methods with and without vector retrieval~\citep{rag,alma,ouyang2025reasoningbank,gepa}.
Detailed baseline implementations are provided in Appendix~\ref{app:baselines}, with implementation notes for every baseline.
